\chapter{FX Spot Microstructure}\label{ch:m2:fx-spot-microstructure}

On 15 December 2010 a client of a large bank's electronic FX platform wrote
to it: ``we have noticed that there were over 300 rejected orders with you
today and the reason is `NACK', could you pls have a look at them''. Two days
later, with no answer, it wrote again: ``We kept receiving top of the book
rates from you and hitting your rate, but we got rejected by you 9 times out of
10''. The New York State Department of Financial Services quoted both messages
five years later, in the consent order by which Barclays paid USD~150 million,
and found no evidence that the bank had ever replied. The rejections came from
a practice called \omterm{def:m2:fx-spot-microstructure:lastlook}{last look}: the bank held each order for a few tens or
hundreds of milliseconds and refused it if the price had moved against the
bank in the meantime. This chapter is about the machinery of electronic FX
spot: how prices are streamed to clients, why the price a client gets depends
on who it is, how \omterm{def:m2:fx-spot-microstructure:lastlook}{last look} works and what it costs
(\cref{fig:m2:fx-spot-microstructure:timeline}), how dealers skew and
internalise, and the code of conduct the market wrote for itself.

\section{Streaming prices and tiered liquidity}

\begin{definition}[Streaming quote, liquidity tier]\label{def:m2:fx-spot-microstructure:stream}
A \emph{streaming quote}\index{streaming quote} is a bid and an ask, for given
sizes, that a liquidity provider publishes continuously to a client or a
platform and updates many times a second; the client trades by sending a
request at a displayed price. A \emph{liquidity tier}\index{liquidity tier} is
the set of spreads and sizes a provider shows to a group of clients, set by its
view of their flow.
\end{definition}

Electronic spot FX is mostly traded on streams. A dealer's pricing engine takes
in the \omterm{def:m2:the-fx-market:venue}{primary venues}, other platforms and its own flow, forms a view of the
mid, and publishes to each client a price around it: tight for clients whose
trades rarely precede price moves, wider for the others, and sometimes wider
again when markets are fast. Principal trading firms stream too, directly to
banks and on platforms (\cref{ch:m2:the-fx-market}). A client with access to
many streams sees the best bid and ask across them, aggregated, and sends its
request to whoever shows the best price.

The dealer's problem is that a stream is a free option. It shows a price for a
few milliseconds, and anyone who knows the market has moved before the stream
has caught up can trade on the stale side. The spread is the dealer's
protection against uninformed noise; tiers, \omterm{def:m2:fx-spot-microstructure:lastlook}{last look} and skews are its
protection against the rest.

\section{Adverse selection}

A market maker's worst clients are those who trade just before the price moves
their way, the adverse selection of One Quant Book 1, chapter 1. In FX spot the
purest form is latency arbitrage: a fast firm sees the price move on one venue
and hits the stale quotes of slower providers elsewhere before they update. A
provider measures its clients with \omterm{def:m2:fx-spot-microstructure:markout}{mark-outs}.

\begin{definition}[Mark-out]\label{def:m2:fx-spot-microstructure:markout}
The \emph{mark-out}\index{mark-out} of a trade at horizon $H$ is the provider's
profit from marking the position at the mid $H$ after the trade: the
half-spread earned minus the move of the mid in the client's favour over $H$.
Averaged over a client's trades at several horizons, from milliseconds to
minutes, it measures how informed the client's flow is.
\end{definition}

A client whose trades have \omterm{def:m2:fx-spot-microstructure:markout}{mark-outs} close to the half-spread at every horizon is
uninformed; its flow is valuable, and it gets the best tier. A client whose
\omterm{def:m2:fx-spot-microstructure:markout}{mark-outs} turn negative within a second is trading on information the provider
did not have; it is shown wider prices, slower confirmations, or nothing.

\section{Last look and hold times}

\begin{definition}[Last look, hold time, reject rate]\label{def:m2:fx-spot-microstructure:lastlook}
\emph{Last look}\index{last look} is a practice by which a liquidity provider
receiving a trade request at its quoted price has a final opportunity, a short
window after the request arrives, to accept or reject it. The \emph{hold
time}\index{hold time} is the length of that window; the \emph{reject
rate}\index{reject rate} is the share of requests rejected. The price check can
be \emph{asymmetric}, rejecting only requests whose price has moved in the
client's favour beyond a threshold, or \emph{symmetric}, rejecting moves beyond
it in either direction.
\end{definition}

The provider's argument for \omterm{def:m2:fx-spot-microstructure:lastlook}{last look} is that it can quote tighter if it may
refuse the trades that arrive on stale prices. The client's objection is that a
rejected request leaves it with the market risk it wanted to shed, at a worse
price, and that an asymmetric check lets the provider keep the trades in which
the price moved its way and refuse the others.

\begin{omfigure}
\centering
\begin{tikzpicture}[font=\small,x=1.3cm]
\draw[thick,-{Stealth}] (0,0) -- (10.4,0) node[right,font=\footnotesize] {time};
\foreach \x/\l in {0/quote sent,2.2/request arrives,6.6/decision} {\draw (\x,0.12) -- (\x,-0.12); \node[below,font=\footnotesize,align=center] at (\x,-0.15) {\l};}
\fill[omThm!15] (2.2,0.2) rectangle (6.6,1.0);
\node[omThm,font=\footnotesize] at (4.4,0.6) {hold time: the provider holds the option};
\draw[omDef,thick,-{Stealth}] (6.6,1.25) -- (7.6,1.85) node[right,font=\footnotesize] {accept: trade at the quoted price};
\draw[omProp,thick,-{Stealth}] (6.6,1.25) -- (7.6,0.65) node[right,font=\footnotesize] {reject: client keeps its risk};
\draw[gray,decorate,decoration={brace,amplitude=4pt,mirror}] (0,-0.9) -- (2.2,-0.9);
\node[gray,font=\footnotesize] at (1.1,-1.3) {latency};
\end{tikzpicture}

\omcaption{The life of a request under \omterm{def:m2:fx-spot-microstructure:lastlook}{last look}. The quote reaches the client,
the client's request reaches the provider after a latency, and the provider
decides at the end of the \omterm{def:m2:fx-spot-microstructure:lastlook}{hold time} whether to trade at the price it quoted.
During the hold the provider, not the client, holds the choice. Schematic.}\label{fig:m2:fx-spot-microstructure:timeline}
\end{omfigure}

\begin{proposition}[What an asymmetric window takes]\label{prop:m2:fx-spot-microstructure:transfer}
Let the price move in the client's favour during the \omterm{def:m2:fx-spot-microstructure:lastlook}{hold time} $h$ by
$X \sim N(0, \sigma^2 h)$, independent of the request, as it is for an
uninformed client. An asymmetric check with threshold $\theta$ rejects when
$X > \theta$, and the client gives up, per request,
\[
\mathbb{E}\bigl[X\,\mathbf 1\{X > \theta\}\bigr] \;=\; \sigma\sqrt{h}\;\varphi\!\left(\frac{\theta}{\sigma\sqrt h}\right).
\]
A symmetric check rejects $|X| > \theta$ and takes nothing in expectation
(\cref{fig:m2:fx-spot-microstructure:tails}).
\end{proposition}

\begin{proof}
For $X = sZ$ with $s = \sigma\sqrt h$, $\mathbb{E}[sZ\mathbf 1\{Z > \theta/s\}] =
s\varphi(\theta/s)$. Under the symmetric check the rejected region is symmetric
about zero, so the rejected moves average to zero.
\end{proof}

\begin{omfigure}
\begin{tikzpicture}
\begin{axis}[name=left,width=6.4cm,height=4.4cm,
  xlabel={move in the client's favour (bp)},yticklabels={},ytick=\empty,
  tick label style={font=\small},label style={font=\small},
  xmin=-0.35,xmax=0.35,ymin=0,ymax=4.4,domain=-0.35:0.35,samples=120,
  title={asymmetric},title style={font=\small}]
\addplot[fill=omProp!35,draw=none,domain=0.05:0.35] {exp(-0.5*(x/0.1)^2)/(0.1*sqrt(2*pi))} \closedcycle;
\addplot[omThm,very thick] {exp(-0.5*(x/0.1)^2)/(0.1*sqrt(2*pi))};
\draw[gray,dashed] (axis cs:0.05,0) -- (axis cs:0.05,4.2);
\node[omProp,font=\scriptsize,align=center] at (axis cs:0.22,2.6) {rejected\\30.9\%};
\end{axis}
\begin{axis}[at={(left.east)},anchor=west,xshift=1.0cm,width=6.4cm,height=4.4cm,
  xlabel={move in the client's favour (bp)},yticklabels={},ytick=\empty,
  tick label style={font=\small},label style={font=\small},
  xmin=-0.35,xmax=0.35,ymin=0,ymax=4.4,domain=-0.35:0.35,samples=120,
  title={symmetric},title style={font=\small}]
\addplot[fill=omProp!35,draw=none,domain=0.05:0.35] {exp(-0.5*(x/0.1)^2)/(0.1*sqrt(2*pi))} \closedcycle;
\addplot[fill=omProp!35,draw=none,domain=-0.35:-0.05] {exp(-0.5*(x/0.1)^2)/(0.1*sqrt(2*pi))} \closedcycle;
\addplot[omThm,very thick] {exp(-0.5*(x/0.1)^2)/(0.1*sqrt(2*pi))};
\draw[gray,dashed] (axis cs:0.05,0) -- (axis cs:0.05,4.2);
\draw[gray,dashed] (axis cs:-0.05,0) -- (axis cs:-0.05,4.2);
\node[omProp,font=\scriptsize,align=center] at (axis cs:0.22,2.6) {rejected\\61.7\%};
\end{axis}
\end{tikzpicture}

\omcaption{The move of the price during a 100-millisecond hold, for an
uninformed request (standard deviation 0.1 basis points), and the requests
rejected with a 0.05-basis-point threshold (shaded). The asymmetric check rejects
only the right tail, whose mean the client loses; the symmetric one rejects both
tails, which cancel. Illustrative.}\label{fig:m2:fx-spot-microstructure:tails}
\end{omfigure}

The symmetric check still costs the client something, the fills it loses, but it
does not select against it. This is the distinction regulators drew. In the New
York order, Barclays had compared the price at the start and at the end of a
hold of tens to hundreds of milliseconds and rejected the orders that had become
unprofitable to it; it made the check
symmetric in September and October 2014, and one platform, with 7\% of its
volume, stayed asymmetric until August 2015.

\begin{omfigure}
\begin{tikzpicture}
\begin{axis}[name=left,width=6.4cm,height=5cm,xmode=log,log basis x=10,
  xlabel={hold time (ms)},ylabel={reject rate (\%)},
  tick label style={font=\small},label style={font=\small},
  xmin=8,xmax=600,ymin=0,ymax=105,grid=major,grid style={gray!25},
  xtick={10,25,50,100,200,500},xticklabels={10,25,50,100,200,500},
  legend columns=2,legend style={font=\footnotesize,at={(1.1,-0.32)},anchor=north,draw=none,fill=none,/tikz/every even column/.append style={column sep=8pt}},
  table/col sep=comma]
\addplot[omProp,very thick,mark=*,mark size=1.2pt] table[x=hold,y=rej_inf_asym]{figdata/markets-2/15-fx-spot-microstructure/hold.csv};
\addplot[omProp,thick,dashed,mark=o,mark size=1.4pt] table[x=hold,y=rej_inf_sym]{figdata/markets-2/15-fx-spot-microstructure/hold.csv};
\addplot[omDef,very thick,mark=*,mark size=1.2pt] table[x=hold,y=rej_uninf_asym]{figdata/markets-2/15-fx-spot-microstructure/hold.csv};
\addplot[omDef,thick,dashed,mark=o,mark size=1.4pt] table[x=hold,y=rej_uninf_sym]{figdata/markets-2/15-fx-spot-microstructure/hold.csv};
\legend{informed asymmetric,informed symmetric,uninformed asymmetric,uninformed symmetric}
\end{axis}
\begin{axis}[at={(left.east)},anchor=west,xshift=1.4cm,width=6.4cm,height=5cm,xmode=log,log basis x=10,
  xlabel={hold time (ms)},ylabel={mark-out (bp)},
  tick label style={font=\small},label style={font=\small},
  xmin=8,xmax=600,ymin=0.35,ymax=0.6,grid=major,grid style={gray!25},
  xtick={10,25,50,100,200,500},xticklabels={10,25,50,100,200,500},
  table/col sep=comma]
\addplot[omThm,very thick,mark=*,mark size=1.2pt] table[x=hold,y=mo_asym]{figdata/markets-2/15-fx-spot-microstructure/hold.csv};
\addplot[omThm,thick,dashed,mark=o,mark size=1.4pt] table[x=hold,y=mo_sym]{figdata/markets-2/15-fx-spot-microstructure/hold.csv};
\draw[gray,dashed] (axis cs:8,0.3785) -- (axis cs:600,0.3785);
\node[font=\scriptsize,gray,anchor=north east] at (axis cs:590,0.3785) {no last look};
\end{axis}
\end{tikzpicture}

\omcaption{A simulated EURUSD stream with 10\% informed requests. Left: \omterm{def:m2:fx-spot-microstructure:lastlook}{reject
rates} by \omterm{def:m2:fx-spot-microstructure:lastlook}{hold time} for informed and uninformed clients under an asymmetric and
a symmetric check with the same 0.05 basis-point threshold. Right: the
provider's one-second \omterm{def:m2:fx-spot-microstructure:markout}{mark-out} per filled request (solid asymmetric, dashed
symmetric). Both checks stop most informed requests; only the asymmetric one
turns longer holds into more profit, taken from uninformed clients.
Illustrative parameters; data: the chapter's tutorial.}\label{fig:m2:fx-spot-microstructure:hold}
\end{omfigure}

\Cref{fig:m2:fx-spot-microstructure:hold} shows both effects on a simulated
stream. Informed requests, which arrive after a move of 0.2 basis points the
quote has not caught up with, are almost all rejected by either check at short
holds. Uninformed requests are rejected more by the symmetric check, which
refuses moves either way, but that costs them nothing on average; the asymmetric
check rejects fewer of them and its \omterm{def:m2:fx-spot-microstructure:markout}{mark-out} rises with the hold, because every
extra millisecond gives the price more time to move in the client's favour and
be refused.

\begin{example}[The cost of a window]\label{ex:m2:fx-spot-microstructure:window}
With the stream's volatility, 0.1 basis points over 100 milliseconds, an
asymmetric check with a threshold of 0.05 basis points and a 100-millisecond
hold takes $0.1\,\varphi(0.5) = 0.0352$ basis points per request from uninformed
clients; with no threshold, 0.0399. On USD~2 billion of requests a day this is
about USD~7\,041 a day, USD~1.76 million over 250 trading days. The symmetric
check with the same threshold rejects 62\% of the requests and takes nothing on
average.
\end{example}

\section{Skewing and internalisation}

\begin{definition}[Quote skewing]\label{def:m2:fx-spot-microstructure:skew}
\emph{Quote skewing}\index{quote skewing} is moving a stream's bid and ask
together, off the provider's estimate of the mid, to attract trades that reduce
its inventory: a provider long euros lowers both prices so that clients buy
euros from it and fewer sell to it (\cref{fig:m2:fx-spot-microstructure:skew}).
\end{definition}

A large provider that receives thousands of client requests a minute can
internalise much of its risk (One Quant Book 1, chapter 10): a sale by one
client is offset by a purchase by another, and only the residual needs to be
hedged on a venue. Skewing makes the netting happen sooner, at the cost of a
small concession on price, and every hedge the provider does not send to a venue
is information it does not reveal. The inventory models behind skews, which
trade off expected spread against the risk of holding a position, are those of
One Quant Book 1 and of the market-making games of
\cref{ch:m2:market-making-games}.

\begin{omfigure}
\centering
\begin{tikzpicture}[font=\small,y=1.1cm]
\draw[gray,dashed] (-0.3,0) node[left,font=\footnotesize] {provider's mid} -- (9.6,0);
\foreach \x/\t/\s in {1.5/flat inventory/0,5/long EUR~20 million/-0.35,8.5/short EUR~20 million/0.35} {
  \draw[omDef,very thick] (\x-0.6,{\s+0.5}) -- (\x+0.6,{\s+0.5}) node[right,font=\footnotesize] {ask};
  \draw[omThm,very thick] (\x-0.6,{\s-0.5}) -- (\x+0.6,{\s-0.5}) node[right,font=\footnotesize] {bid};
  \node[font=\footnotesize,align=center] at (\x,-1.45) {\t};
}
\draw[omProp,thick,-{Stealth}] (5,-0.05) -- (5,-0.32);
\draw[omProp,thick,-{Stealth}] (8.5,0.05) -- (8.5,0.32);
\end{tikzpicture}

\omcaption{\omterm{def:m2:fx-spot-microstructure:skew}{Quote skewing}. With no position the provider quotes symmetrically
around its mid; long euros, it lowers both prices, so that buyers come to it and
sellers go elsewhere; short, it raises them. The spread is unchanged, the
centre moves. Schematic.}\label{fig:m2:fx-spot-microstructure:skew}
\end{omfigure}

\section{The global code of conduct}

\begin{definition}[FX Global Code]\label{def:m2:fx-spot-microstructure:code}
The \emph{FX Global Code}\index{FX Global Code} is a set of principles of good
practice for the wholesale FX market, maintained by the Global Foreign Exchange
Committee of central banks and private-sector participants. It is not law:
participants adhere by signing a statement of commitment.
\end{definition}

The Code arrived in 2017, two years after the enforcement cases over the
fixing of \cref{ch:m2:fixings-and-flows}, and \omterm{def:m2:fx-spot-microstructure:lastlook}{last look} is among its most
argued-over sections. Principle~17, as updated in
December 2024, requires providers that use \omterm{def:m2:fx-spot-microstructure:lastlook}{last look} to be transparent about it:
to say whether and how price changes in either direction affect the decision,
how long it takes and why they use it. It defines \omterm{def:m2:fx-spot-microstructure:lastlook}{last look} as a risk control, to
check that a request is valid and that its price is still consistent with the
market, and forbids using it to gather information or to trade on the request
during the window, whether by adjusting prices or by hedging. In August 2021
the Committee published templates on which providers disclose their practices.

\begin{remark}[What the Code leaves open]\label{rem:m2:fx-spot-microstructure:open}
The Code does not ban asymmetric checks, nor does it set \omterm{def:m2:fx-spot-microstructure:lastlook}{hold times}. It asks
providers to disclose them and clients to read the disclosures; the pressure
that remains is competitive. Clients measure their fill ratios, \omterm{def:m2:fx-spot-microstructure:lastlook}{hold times} and
\omterm{def:m2:fx-spot-microstructure:markout}{mark-outs} by provider, as the tutorial does, and route their requests
accordingly.
\end{remark}

\begin{dated}{2026-09}{The Code today}\label{dat:m2:fx-spot-microstructure:code}
The \omterm{def:m2:fx-spot-microstructure:code}{FX Global Code}'s current text is dated December 2024; its Principle~17 on
\omterm{def:m2:fx-spot-microstructure:lastlook}{last look} is as summarised above. The Global Foreign Exchange Committee, set up
in May 2017 to maintain it, published in August 2021 a guidance paper on \omterm{def:m2:fx-spot-microstructure:lastlook}{last
look} and standard templates on which liquidity providers disclose their
practices. Adherence remains voluntary, by statement of commitment.
\end{dated}

\section{Tutorial: last-look transaction-cost analysis}\label{tut:m2:fx-spot-microstructure:tca}

\begin{tutorial}
\textbf{Goal.} Simulate a stream with informed and uninformed clients, apply
asymmetric and symmetric last-look checks, and measure \omterm{def:m2:fx-spot-microstructure:lastlook}{reject rates}, fill ratios
and \omterm{def:m2:fx-spot-microstructure:markout}{mark-outs} as a function of the \omterm{def:m2:fx-spot-microstructure:lastlook}{hold time}. \textbf{End state:}
\cref{fig:m2:fx-spot-microstructure:hold}, \cref{ex:m2:fx-spot-microstructure:window}
and the numbers of the weekend problem.
\begin{tutsteps}
\item \textbf{The check and the stream}: each request's move during the hold
and after it.
\omcode{code/firm/lastlook/firm_lastlook.py}{24}{43}{The last-look check and a simulated stream of requests.}\label{lst:m2:fx-spot-microstructure:sim}
\item \textbf{The analysis}: fill ratio, \omterm{def:m2:fx-spot-microstructure:lastlook}{reject rates} by client type, \omterm{def:m2:fx-spot-microstructure:markout}{mark-out},
and the closed form of \cref{prop:m2:fx-spot-microstructure:transfer}.
\omcode{code/firm/lastlook/firm_lastlook.py}{46}{69}{Transaction-cost analysis and the expected transfer.}\label{lst:m2:fx-spot-microstructure:tca}
\item \textbf{Run} \texttt{lastlook\_demo.by\_hold} for each policy,
\texttt{lastlook\_demo.window()} and \texttt{fig\_lastlook.py}.
\end{tutsteps}
\textbf{What to change next.} Replace the informed clients' fixed edge by a
random one, and find the \omterm{def:m2:fx-spot-microstructure:lastlook}{hold time} at which a symmetric check has rejected 90\%
of them; then plot the uninformed clients' fill ratio at that hold.
\end{tutorial}

\section{Build: the last-look analyser}\label{bld:m2:fx-spot-microstructure:lastlook}

\begin{build}
\bfield{Purpose} The miniature firm both takes and gives liquidity in FX spot:
as a taker it must know what each provider's \omterm{def:m2:fx-spot-microstructure:lastlook}{last look} costs it; as a maker it
must be able to show its clients that its own check is symmetric and short.
\bfield{Interface} \texttt{check(move, threshold, policy)};
\texttt{simulate(n, sigma, hold, threshold, policy, informed\_share, edge,
horizon, seed)} returning \texttt{Request} records; \texttt{tca(requests,
half)}; \texttt{expected\_transfer(sigma, hold, threshold, policy)}.
\bfield{Rules} Moves in basis points, positive in the client's favour;
Brownian mid; policies \texttt{none}, \texttt{asymmetric}, \texttt{symmetric};
seeded simulation.
\bfield{Acceptance tests} \texttt{code/firm/lastlook/tests/}: the policies'
definitions; no \omterm{def:m2:fx-spot-microstructure:lastlook}{last look} fills everything; the simulated asymmetric transfer
within 2\% of the closed form; informed requests rejected far more often than
uninformed ones.
\bfield{Stretch} Read real fill and reject records (a FIX log) and compute each
provider's hold-time distribution and \omterm{def:m2:fx-spot-microstructure:markout}{mark-outs}; test statistically whether a
provider's rejects depend on the direction of the move.
\end{build}

\begin{omsources}
\item New York State Department of Financial Services, Consent Order, In the
Matter of Barclays Bank PLC, November 2015.
\item Global Foreign Exchange Committee, \emph{FX Global Code}, updated
December 2024; press release on last-look guidance and disclosure templates,
August 2021.
\item Bank for International Settlements, \emph{Quarterly Review}, December
2022, box on the fragmented spot market.
\end{omsources}

\section{Exercises}

\begin{exercise}[$\star$]\label{exo:m2:fx-spot-microstructure:1}
A provider sells a client EUR~1 million at 1.1464 when the mid is 1.1463. One
second later the mid is 1.1465. Give the provider's one-second \omterm{def:m2:fx-spot-microstructure:markout}{mark-out} in
dollars and in basis points.
\end{exercise}

\begin{exercise}[$\star$]\label{exo:m2:fx-spot-microstructure:2}
Give the standard deviation of the price move over holds of 25, 100 and 400
milliseconds at 0.01 basis points per square-root millisecond.
\end{exercise}

\begin{exercise}[$\star$]\label{exo:m2:fx-spot-microstructure:3}
What must a provider that uses \omterm{def:m2:fx-spot-microstructure:lastlook}{last look} disclose under Principle~17, and what
must it not do during the window?
\end{exercise}

\begin{exercise}[$\star\star$]\label{exo:m2:fx-spot-microstructure:4}
Give the transfer per request of an asymmetric check at a 25-millisecond hold
and zero threshold, and compare it with the 100-millisecond hold.
\end{exercise}

\begin{exercise}[$\star\star$]\label{exo:m2:fx-spot-microstructure:5}
A provider long EUR~20 million skews its EURUSD stream by 0.1 pip. Which way,
and why may this be cheaper than hedging on a \omterm{def:m2:the-fx-market:venue}{primary venue}?
\end{exercise}

\begin{exercise}[$\star\star$]\label{exo:m2:fx-spot-microstructure:6}
Why does a symmetric check reject more uninformed requests than an asymmetric
one with the same threshold, yet cost them less?
\end{exercise}

\begin{exercise}[$\star\star\star$]\label{exo:m2:fx-spot-microstructure:7}
\emph{Coding.} From \texttt{by\_hold}, give the informed and uninformed \omterm{def:m2:fx-spot-microstructure:lastlook}{reject
rates} at 100 milliseconds under each policy, and the \omterm{def:m2:fx-spot-microstructure:markout}{mark-outs}.
\end{exercise}

\begin{exercise}[$\star\star\star$]\label{exo:m2:fx-spot-microstructure:8}
\emph{Find the flaw.} ``Our \omterm{def:m2:fx-spot-microstructure:lastlook}{reject rate} is only 5\%, so our \omterm{def:m2:fx-spot-microstructure:lastlook}{last look} costs
our clients nothing.'' Correct it.
\end{exercise}

\section{Problem: The Asymmetric Window}

\begin{problem}[{Weekend problem --- what a last-look window takes from a client}]\label{pb:m2:fx-spot-microstructure:1}
An asset manager sends USD~2 billion of EURUSD requests a day to a provider
that uses \omterm{def:m2:fx-spot-microstructure:lastlook}{last look} with a 100-millisecond hold and a 0.05-basis-point
threshold. The mid moves 0.01 basis points per square-root millisecond, the
half-spread is 0.4 basis points, and the manager's flow is uninformed.

\medskip\noindent\textbf{Part I --- The window.}
\begin{enumerate}
\item Give the standard deviation of the move over the hold.
\item Under an asymmetric check, what share of requests is rejected?
\item Under a symmetric check with the same threshold?
\item Give the transfer per request under each.
\item Give the asymmetric transfer with no threshold.
\end{enumerate}

\medskip\noindent\textbf{Part II --- In money.}
\begin{enumerate}[resume]
\item Give the daily cost of the asymmetric window.
\item Give it over 250 trading days.
\item How does it compare with the half-spread the manager pays?
\item Give the cost with a 25-millisecond hold and no threshold.
\item Why does the cost grow with the square root of the hold?
\end{enumerate}

\medskip\noindent\textbf{Part III --- What the manager sees.}
\begin{enumerate}[resume]
\item Which statistics reveal an asymmetric check?
\item Why is a low \omterm{def:m2:fx-spot-microstructure:lastlook}{reject rate} not enough?
\item What should the manager ask the provider, citing the Code?
\item What else does a rejection cost, beyond the transfer?
\item Why might the manager still prefer a provider with \omterm{def:m2:fx-spot-microstructure:lastlook}{last look}?
\end{enumerate}

\medskip\noindent\textbf{Part IV --- Judgement.}
\begin{enumerate}[resume]
\item What did the New York order find wrong with Barclays's check?
\item Is a symmetric check fair to the client?
\item How does a firm-price, no-last-look stream protect itself instead?
\item State the \emph{named result}: what the asymmetric window takes from the
manager each year.
\item In one sentence: what does a last-look window cost a client?
\end{enumerate}
\end{problem}

\section{Interview questions}

\begin{interviewq}[$\star$ \iqroles{trader, developer}]\label{iq:m2:fx-spot-microstructure:1}
What is \omterm{def:m2:fx-spot-microstructure:lastlook}{last look}, and why do liquidity providers use it?
\end{interviewq}

\begin{interviewq}[$\star$ \iqroles{researcher, trader}]\label{iq:m2:fx-spot-microstructure:2}
How would you tell, from your fills, whether a provider's \omterm{def:m2:fx-spot-microstructure:lastlook}{last look} is
asymmetric?
\end{interviewq}

\begin{interviewq}[$\star\star$ \iqroles{researcher}]\label{iq:m2:fx-spot-microstructure:3}
What is a \omterm{def:m2:fx-spot-microstructure:markout}{mark-out}, and how do you use \omterm{def:m2:fx-spot-microstructure:markout}{mark-outs} to tier clients?
\end{interviewq}

\begin{interviewq}[$\star\star$ \iqroles{trader}]\label{iq:m2:fx-spot-microstructure:4}
You are long EUR~50 million from client flow. Do you skew, hedge, or wait?
\end{interviewq}

\begin{interviewq}[$\star\star$ \iqroles{bank, trader}]\label{iq:m2:fx-spot-microstructure:5}
What does the \omterm{def:m2:fx-spot-microstructure:code}{FX Global Code} require of \omterm{def:m2:fx-spot-microstructure:lastlook}{last look}, and what does it not
require?
\end{interviewq}

\begin{interviewq}[$\star\star\star$ \iqroles{developer}]\label{iq:m2:fx-spot-microstructure:6}
Design a pricing engine that streams EURUSD to a thousand clients in several
tiers and applies a symmetric last-look check, within a latency budget of tens
of microseconds.
\end{interviewq}
